% ================================================================
% ==== FILE: content/cycles/cycles_k6.tex
% ================================================================

\subsubsection{Zyklen auf \texorpdfstring{$K_6$}{K_6}}
\label{sec:cycles-k6}

Das Niveau$K_6$ist das kognitive Kontinuum: ein System mit Repräsentation
Achsen, prädiktive Dynamik, Bindungsmechanismen, Gedächtnisintegration und
Strukturierte interne Modelle.
Zyklen weiter$K_6$entsprechen wiederkehrenden kognitiven Prozessen, die sich stabilisieren
Darstellungen, bewahren die Kohärenz und ermöglichen Vorhersagen und Lernen.

Die Struktur der Zyklen weiter$K_6$wird durch die kognitiven Achsen bestimmt
\[
 A_{\mathrm{rep}},\quad
 A_{\mathrm{bind}},\quad
 A_{\mathrm{pred}},\quad
 A_{\mathrm{mem}},\quad
 A_{\mathrm{attn}},
\]
die Potenziale
\[
 P_{\mathrm{rep}},\, P_{\mathrm{pred}},\, P_{\mathrm{error}},\, 
 P_{\mathrm{salience}},\, P_{\mathrm{stability}},
\]
flows $J_6$ (the corresponding derivation note), and thresholds:
\[
 \Theta_{\mathrm{PE}},\quad
 \Theta_{\mathrm{bind}},\quad
 \Theta_{\mathrm{coh}},\quad
 \Theta_{\mathrm{expressivity}},\quad
 \Theta_{\mathrm{noise}}.
\]

\subsubsection{Überblick über Zyklen auf Erkenntnisebene}

Kognitive Zyklen entstehen, weil kognitive Zustände aktiv aufrechterhalten werden müssen
gegen Lärm, Instabilität und Repräsentationsverfall.
Jeder Zyklus läuft$K_6$drückt ein geschlossenes Regelkreisregime aus:

\[
  \mathrm{perception}
  \;\to\;
  \mathrm{prediction}
  \;\to\;
  \mathrm{comparison}
  \;\to\;
  \mathrm{error\;evaluation}
  \;\to\;
  \mathrm{update}
  \;\to\;
  \mathrm{re-stabilisation}.
\]

Diese Zyklen bilden die strukturelle Grundlage der Erkenntnis.

\subsubsection{1. Vorhersage-Verarbeitungszyklus$C_{\mathrm{PE}}$}

Der zentrale kognitive Zyklus entsteht aus der prädiktiven Verarbeitung, gesteuert durch
the flow $J_{\mathrm{PE}}$ (the corresponding derivation note):

\[
 C_{\mathrm{PE}} = 
 \left\{
   (P_{\mathrm{pred}}(t),
    P_{\mathrm{rep}}(t),
    P_{\mathrm{error}}(t))
 \right\}_{t_0}^{t_0+\tau_{\mathrm{PE}}}.
\]

Der Zyklus existiert, wenn die Vorhersage-Fehler-Schleife stabil ist:

\[
 P_{\mathrm{error}}(t+\tau_{\mathrm{PE}}) = P_{\mathrm{error}}(t),
\]

und kognitive Spannung befriedigt:

\[
 T_{\mathrm{PE}} < \Theta_{\mathrm{PE}}.
\]

Violation of $\Theta_{\mathrm{PE}}$ generates cognitive collapse or chaotic
Vorhersagefehler.

\subsubsection{2. Bindungszyklus$C_{\mathrm{bind}}$}

Kognitive Darstellungen müssen ihre Komponenten miteinander verbinden
representational axes $A_{\mathrm{bind}}$.

Der Bindungszyklus:

\[
 \mathrm{feature\ extraction}
 \to
 \mathrm{coherence\ formation}
 \to
 \mathrm{integration}
 \to
 \mathrm{comparison}
 \to
 \mathrm{re-binding}.
\]

Formal:
\[
 C_{\mathrm{bind}} =
 \left\{
   A_{\mathrm{bind}}(t), P_{\mathrm{rep}}(t)
 \right\}_{t_0}^{t_0+\tau_{\mathrm{bind}}}.
\]

Existenz erfordert:
\[
 T_{\mathrm{bind}} < \Theta_{\mathrm{bind}}.
\]

Dieser Zyklus entspricht der wahrnehmungsbezogenen und konzeptionellen Kohärenz.

\subsubsection{3. Aufmerksamkeits-Salienz-Zyklus$C_{\mathrm{attn}}$}

Attention modulates the salience potential $P_{\mathrm{salience}}$ and
Leitet Flüsse selektiv weiter$J_6$.

Der Zyklus ist:

\[
  \mathrm{salience\ detection}
  \to
  \mathrm{attentional\ shift}
  \to
  \mathrm{enhancement}
  \to
  \mathrm{integration}
  \to
  \mathrm{decay}.
\]

Die Wiederholungsbedingung:

\[
  P_{\mathrm{salience}}(t+\tau_{\mathrm{attn}}) = P_{\mathrm{salience}}(t).
\]

Stabilität erfordert:
\[
  T_{\mathrm{attn}} < \Theta_{\mathrm{coh}}.
\]

\subsubsection{4. Arbeitsspeicherzyklus$C_{\mathrm{wm}}$}

Das Arbeitsgedächtnis beinhaltet die wiederkehrende Aufrechterhaltung repräsentativer Zustände.
Die klassische Schleife:

\[
 \mathrm{load}
 \to
 \mathrm{maintain}
 \to
 \mathrm{transform}
 \to
 \mathrm{offload}.
\]

Formal:
\[
 C_{\mathrm{wm}} = 
 \left\{
   P_{\mathrm{mem}}(t), A_{\mathrm{mem}}(t)
 \right\}_{t_0}^{t_0+\tau_{\mathrm{wm}}}.
\]

Existenz erfordert:
\[
 T_{\mathrm{wm}} < \Theta_{\mathrm{expressivity}}.
\]

Diese Schwelle verkörpert die minimale Ausdrucksfähigkeit von$K_6$.

\subsubsection{5. Konzeptioneller Zyklus$C_{\mathrm{concept}}$}

Ein Konzept wird durch Wiederholung über Transformationsachsen hinweg stabilisiert:

\[
  C_{\mathrm{concept}}
    = \left\{
      A_{\mathrm{rep}}(t),
      P_{\mathrm{rep}}(t),
      P_{\mathrm{stability}}(t)
    \right\}.
\]

Es ist definiert durch:

\[
  P_{\mathrm{stability}}(t+\tau) = P_{\mathrm{stability}}(t),
\]

und Kohärenzschwelle:

\[
  T_{\mathrm{rep}} < \Theta_{\mathrm{coh}}.
\]

Dieser Zyklus verallgemeinert symbolische Stabilität, ohne dass Sprache erforderlich ist.

\subsubsection{6. Kognitiver Rhythmuszyklus$C_{\mathrm{cog\text{-}rhythm}}$}

Die Wahrnehmung beruht oft auf der rhythmischen Koordination (Vorläufer der
neuronale Rhythmen von$K_5'$aber jetzt intern zur Darstellung).

Der kognitive Rhythmuszyklus:

\[
 P_{\mathrm{rep}}(t) = P_0 + A\sin(\omega t),\qquad
 A < \Theta_{\mathrm{noise}}.
\]

Dieser Zyklus sorgt für die Kohärenz der Darstellungsrahmen.

\subsubsection{7. Multiskalen-Integrationszyklus$C_{\mathrm{integrate}}$}

Kognition integriert sich über Zeitskalen hinweg:

\[
 \mathrm{fast\ prediction}
 \leftrightarrows
 \mathrm{slow\ model\ update}.
\]

Formal ein langsam-schneller Zyklus:
\[
 C_{\mathrm{integrate}} =
 \left\{
   (P_{\mathrm{pred}}^{\mathrm{fast}}(t),
    P_{\mathrm{rep}}^{\mathrm{slow}}(t))
 \right\}.
\]

Der Zyklus ist nur dann lebensfähig, wenn:

\[
 T_{\mathrm{multi}} < \min(\Theta_{\mathrm{PE}},\Theta_{\mathrm{expressivity}}).
\]

\subsubsection{8. Metakognitiver Zyklus$C_{\mathrm{meta}}$}

Metakognition entspricht Zyklen, in denen:
\[
 \mathrm{model}
 \to 
 \mathrm{self-monitoring}
 \to
 \mathrm{evaluation}
 \to
 \mathrm{correction}
 \to
 \mathrm{model}.
\]

Dies ist der früheste Vorläufer der metatheoretischen Kontinua$K_9$.

Stabilitätsbedingung:
\[
 T_{\mathrm{meta}} < \Theta_{\mathrm{coh}}.
\]

\subsubsection{Metriken kognitiver Zyklen}

Using the scoped cycle metrics (the corresponding derivation note):

\paragraph{Length.}
\[
  L(C)=\oint d\Omega(K_6).
\]

\paragraph{Efficiency.}
\[
  C_{\mathrm{eff}}
    = \frac{\oint J_6\cdot dA}{L(C)}.
\]

\paragraph{Stability.}
\[
SC)
  = \min_{t\in C}
    d(\Omega(K_6),\partial\Omega(K_6)).
\]

\paragraph{Weight.}
\[
WC)
 = F(C_{\mathrm{eff}}, S(C), P_{\mathrm{pred}},
      P_{\mathrm{mem}}, P_{\mathrm{error}}).
\]

\subsubsection{Zusammenbruch von \texorpdfstring{$K_6$}{K_6} Zyklen}

Cognitive collapse (the corresponding derivation note) occurs when:
\[
  T_{\mathrm{PE}}>\Theta_{\mathrm{PE}},
  \quad
  T_{\mathrm{bind}}>\Theta_{\mathrm{bind}},
  \quad
  T_{\mathrm{coh}}>\Theta_{\mathrm{coh}},
  \quad
  T_{\mathrm{expressivity}}>\Theta_{\mathrm{expressivity}}.
\]

Dann:
\[
 C_j \to \varnothing,
 \qquad
 \Omega(K_6)=\varnothing,
 \qquad
 k_6\to 0.
\]

\subsubsection{Kontinuität von \texorpdfstring{$K_5$}{K_5} zu \texorpdfstring{$K_6$}{K_6}}

The transition $K_5' \to K_6$ (the corresponding derivation note) corresponds to the rise
von:
\[
 A_{\mathrm{rep}},\;
 A_{\mathrm{bind}},\;
 A_{\mathrm{pred}},\;
 A_{\mathrm{mem}},
\]
zusammen mit neuen Schwellenwerten:

\[
 \Theta_{\mathrm{PE}},\;
 \Theta_{\mathrm{bind}},\;
 \Theta_{\mathrm{coh}},\;
 \Theta_{\mathrm{expressivity}}.
\]

Dies ist ein Dimensionsübergang, der durch universelle Bedingungen bestimmt wird:
\[
  T > \Theta_{\mathrm{dim}}.
\]

Die entstehenden Zyklen gehen weiter$K_6$Vervollständigen Sie die Bildung eines vollständig kognitiven Kontinuums.
